\section{Experimental Setup}

\subsection{Gold Standard}
The gold standard was manually constructed with 1,582 instances, consisting of
English words alongside their Tagalog phonetically-spelled versions. The words
were separated into 23 grapheme groups (e.g., ed, ough, tch). These grapheme groups
were based on a phoneme-to-grapheme chart, where digraphs up to quadgraphs were
selected to represent as many words as possible that contain corresponding
grapheme sequences. Words were collected by searching each grapheme on
Merriam-Webster's Word Finder using the ``Containing in order'' filter and selecting
the most commonly used words.

The phonetically-spelled versions did not include the letters \emph{z}, \emph{x}, 
\emph{c}, and \emph{v}, as those letters are not present in traditional Tagalog.
Each word was constructed based on the Basilect pronunciations of the English word
and was validated using Inter-Annotator Agreement. For reliability, three Filipino
annotators compared their own versions of each word and selected the one that all
annotators agreed fit best.

\subsection{Metrics}
Metrics were measured through \textbf{Word Error Rate (WER)} and \textbf{Character Error
Rate (CER)}. WER measured the error on a whole-word basis (comparing the actual output to
our expected output in the dataset), while CER measured the error on a per-character basis.
